% Generated draft: requires mathematical and source review.
\documentclass[11pt]{article}
\usepackage[margin=25mm]{geometry}
\usepackage{fontspec}
\setmainfont{Latin Modern Roman}
\usepackage{amsmath,amssymb,mathtools,unicode-math}
\setmathfont{Latin Modern Math}
\usepackage{enumitem,needspace,xurl,hyperref}
\setlength{\parindent}{0pt}
\setlength{\parskip}{5pt}
\setlength{\emergencystretch}{3em}
\newcommand{\N}{\mathbb{N}}
\newcommand{\Z}{\mathbb{Z}}
\newcommand{\Q}{\mathbb{Q}}
\newcommand{\R}{\mathbb{R}}
\newcommand{\C}{\mathbb{C}}
\newcommand{\F}{\mathbb{F}}
\newcommand{\A}{\mathbb{A}}
\newcommand{\PP}{\mathbb{P}}
\newcommand{\E}{\mathbb{E}}
\newcommand{\eps}{\varepsilon}
\newcommand{\dd}{\,\mathrm{d}}
\begin{document}
\section*{UnsolvedMath ID: 11000308}
\textbf{Source problem code:} catalogue entry
\subsection{Definitions and mathematical statement}
Let \(\mathcal K_g\) be the source's Johnson kernel and let \(d_1\) be its previously defined class/homomorphism. The main task is to give an explicit method calculating \(d_1(\varphi)\) for every \(\varphi\in\mathcal K_g\). Its definition, normalization, codomain and input encoding are absent, so no numerical formula for a separating twist or other element is fabricated. A word-input formulation must specify generators and promise membership in \(\mathcal K_g\).

The first concrete question concerns the Magnus representation
\[r:\mathcal I_{g,1}\to GL(2g;\Z[H]),\]
where \(H\) is the relevant first-homology module and \(\Z[H]\) its group ring. Does it detect \(d_1\)? Once the domain/lift of \(\mathcal K_g\) is identified, detection means that equal representation values give equal \(d_1\)-values; for a homomorphism \(d_1\) this is equivalent to
\[\ker r\cap\mathcal K_g\subseteq\ker d_1,\]
or factorization of \(d_1\) through \(r(\mathcal K_g)\). The source cites Suzuki [103] for nonfaithfulness of \(r\), but explicitly only suggests that this inclusion might fail. Nonfaithfulness alone does not decide it. Since the displayed representation is on \(\mathcal I_{g,1}\) while \(\mathcal K_g\) may use the closed-surface convention, compatible lifts/quotients are essential missing data.

The second concrete question uses Jones's representation [45] of the hyperelliptic mapping class group \(\Delta_g\): does its restriction to \(\Delta_g\cap\mathcal K_g\) detect \(d_1\)? With that representation denoted by \(J\), the homomorphism reading asks
\[\ker J\cap\mathcal K_g\subseteq\ker d_1\]
inside the common domain, equivalently whether the restricted \(d_1\) factors through its image. Its coefficient/parameter/normalization conventions are not supplied. Kasahara [46] is cited only as related work for \(g=2\).

The source also names Endo [14] and Morifuji [71] for hyperelliptic work, Kitano [54], and Hain--Reed [33] as related to \(d_1\). It reports a Biss--Farb [5] claim that \(\mathcal K_g\) is not finitely generated for all \(g\geq3\), Mess [69] for \(\mathcal K_2\) as an infinitely generated free group, and says finite generation of its abelianization was still unknown, referring to Problem 2.2 of [80]. The historical Biss--Farb statement must be checked against primary sources and subsequent corrections before being treated as a current theorem; even non-finite-generation alone would not settle abelianization. Finally it reports Kawazumi [50] developing harmonic Magnus expansions giving differential forms representing Mumford--Morita--Miller classes on the universal curve over \(\mathcal M_g\). This contextual report is not an explicit solution to the \(d_1\) calculation or detection questions.
\subsection{Short English statement}
Compute the specified \(d_1\) invariant on every Johnson-kernel element and determine whether Magnus or hyperelliptic Jones representations retain it. The invariant definition, compatible domains and several historical status assertions need checking.
\subsection{Sources}
\begin{itemize}
\item Local UnsolvedMath registry, permanent ID \(11000308\), exact source code \texttt{catalogue entry}.
\item Farb, \emph{Problems on Mapping Class Groups and Related Topics}, Problem 3.5, PDF page 363: \url{https://www.math.uchicago.edu/~farb/papers/mcgbook.pdf}. Registry access 2026-07-29; pdf-font-structure extraction.
\item Suzuki [103]; Jones [45]; Kasahara [46]; Endo [14]; Morifuji [71]; Kitano [54]; Hain--Reed [33]; Biss--Farb [5]; Mess [69]; Problem 2.2 of [80]; and Kawazumi [50]. All are supplied internal references without full bibliographic entries.
\end{itemize}
\end{document}
